\section{工程实现与代码结构}

以下从目录组织、关键组件以及训练评测流程三个方面，介绍本设计的工程实现细节。

\subsection{项目目录结构}

代码库采用功能模块化的方式进行组织，各顶层目录的职责如下：

\begin{lstlisting}[language=bash,caption={项目目录结构}]
project_root/
  src/
    data/          # 数据加载、预处理、CLIP特征缓存
    models/        # 模型定义(CLIP, A2 encoder, fusion, caption model)
    train/         # 训练脚本(baseline, fusion, alignment等)
    eval/          # 评估(generation, retrieval, sanity check, metrics)
    losses/        # 损失函数(contrastive, EEG augmentation, similarity)
    utils/         # 工具(config loader, seed, checkpoint, logger)
  scripts/         # 启动脚本和数据工具
  configs/         # YAML配置文件
  outputs/         # 训练输出(checkpoint, metrics, report)
  tests/           # 单元测试
\end{lstlisting}

\subsection{核心模块说明}

表~\ref{tab:code-modules} 列出了主要模块及其功能。

\begin{table}[htbp]
  \centering
  \caption{核心代码模块}
  \label{tab:code-modules}
  \begin{tabularx}{\linewidth}{l X}
    \toprule
    模块 & 功能 \\
    \midrule
    \texttt{src/data/dataset.py} & EEGVisionCaptionDataset，支持 JSONL manifest 的数据加载器，兼容 Thought2Text、EEG-ImageNet、THINGS-EEG2 等格式 \\
    \texttt{src/data/clip\_cache.py} & CLIP 视觉特征预缓存模块 \\
    \texttt{src/models/eeg\_encoder.py} & 13 种 EEG encoder variant 定义及 \texttt{build\_eeg\_encoder} 工厂函数 \\
    \texttt{src/models/fusion.py} & GatedFusion 门控残差融合模块 \\
    \texttt{src/models/caption\_model.py} & SoftPromptCaptionModel 与内置 byte-level 回退 \\
    \texttt{src/train/train\_semantic\_fusion.py} & A2 语义融合训练脚本 \\
    \texttt{src/train/train\_align.py} & EEG-to-CLIP 对齐训练（InfoNCE+MSE）\\
    \texttt{src/eval/generate.py} & 多模式生成式评估 \\
    \texttt{src/eval/semantic\_fusion\_classifier\_eval.py} & A2 语义融合分类评估 \\
    \bottomrule
  \end{tabularx}
\end{table}

\subsection{训练与评估流程}

\textbf{训练流程}可分为六个步骤。(1) 数据准备阶段：编写 JSONL manifest 并预计算 CLIP 图像特征存入缓存。(2) A2 EEG 编码器对齐（按需执行）：以 InfoNCE 对比损失为目标，将 EEG 表示映射至 CLIP 语义空间。(3) A2 语义融合训练：对融合模块与分类头部进行联合优化。(4) VTF token 级融合训练。(5) 生成式模型微调：采用 Qwen2-VL 结合 LoRA 低秩适配。(6) 全量评测：遍历全部退化条件与 EEG 模式的笛卡尔积。

\textbf{输出产物}：每一轮训练均在 \texttt{outputs/<run\_name>/} 目录下保存最优 checkpoint、运行日志、metrics JSON 以及 report。生成式评测结果以 JSONL 格式存档，每行包含 image\_id、mode、reference、prediction 字段。

\subsection{遇到的工程问题}

开发过程中遇到并处理了以下几项工程问题：(1) 增强后的视觉 token 维度与 Qwen2-VL 的输入规格存在差异，需在 VTF 之后接入本设计自行实现的轻量 prefix adapter，完成格式对齐与维度匹配。(2) 部分 BLIP 生成的伪 caption 出现内容重复、文本过长或与类别标签不符的现象，本设计采用最小长度过滤（$\geq$ 3 个词）配合正则表达式去噪的方式进行清理。(3) 全量测试规模为 1997 $\times$ 6 $\times$ 4 $\approx$ 48000 次推理，生成式模型的自回归解码构成主要时间瓶颈，一轮完整评估耗时约 2--4 小时。(4) 主要训练任务全部在单张 48 GB GPU 上完成，未引入 DeepSpeed 或多 GPU 分布式策略。